\scanpage{20}
\hfill\fbox{Lemmes sur coeff. binomiaux.}

\uline{Lemme} 1 Soient $p$ nb premier, $i,j\in\mathbb{N}$, avec \unsure{des développements|du développement} $p$-\unsure{adiques|adique}
\[
i=\sum_{0\leq\varepsilon\leq{e}}u_\varepsilon p^\varepsilon
\qquad j=\sum_{0\leq\varepsilon\leq e}v_\varepsilon p^\varepsilon
\qquad 0\leq u_\varepsilon,v_\varepsilon\leq p-1,
\]
soit
\[
c_{ij}=\frac{(i+j)!}{i!\,j!}\in\mathbb{Z}\qquad\text{(coeff. binômial)}
\]
alors on a
\[
c_{ij}\not\equiv0\ (p)\quad\text{sss}\quad u_\varepsilon+v_\varepsilon\leq p-1\quad\forall\varepsilon\quad(0\leq\varepsilon\leq\unsure{\nu|v})
\]
et \deleted{\unsure{pour}} on a en tous cas
\[
c_{ij}\equiv\prod_{0\leq\varepsilon\leq\nu}c_{u_\varepsilon,v_\varepsilon}\qquad(p)
\]

\margin{\uline{Cor} 1}% slanted
On écrit d'abord $i=u+pr,\quad j=v+ps\qquad 0\leq u,v\leq p-1,\quad r,s\geq0$ et on montre que
\[
\left\{\begin{aligned}
&c_{ij}\not\equiv0\ (p)\quad\text{sss}\quad u+v\leq p-1\quad\text{et}\quad c_{r,s}\not\equiv0\quad(p)\\
&c_{ij}\equiv c_{u,v}\,c_{r,s}\qquad(p)
\end{aligned}\right.\tag*{$(*)$}
\]
en étudiant le développement, dans $\mathbb{F}_p[V,Y]$, de $(X+Y)^{i+j}$. \deleted{\unsure{ss}}On a
\[
i+j=p\underbrace{(r+s)}_{t}+\underbrace{u+v}_{\deleted{\unsure{\lambda}}w}
\]
si $w\leq p-1$ on écrit
\[
\begin{aligned}
(X+Y)^{\unsure{j|i}+j}&=\bigl((X+Y)^p\bigr)^{t\deleted{\unsure{r+s}}}(X+Y)^{w}=(X^p+Y^p)^{t\deleted{\unsure{r+s}}}(X+Y)^w\\
&=\sum_{\substack{0\leq\alpha\leq w\\0\leq\beta\leq t}}\binom{w}{\alpha}\binom{t}{\beta}X^{\alpha+p\beta}Y^{(w-\alpha)+p(t-\beta)},
\end{aligned}
\]
où tous les monômes \unsure{écrit|écrits} dans la somme sont à exposants distincts. On a donc comme coeff. de $X^iY^j$ celui correspondant au terme $i=\alpha+p\beta$ i.e.
\[
\alpha=u,\quad\beta=r\qquad\text{i.e.}
\]
\[
c_{i,j}\equiv\binom{w}{u}\binom{t}{r}=c_{u,v}\,c_{r,s}\qquad(p)
\]
et il est immédiat que $c_{u,v}\not\equiv0\ (p)$, car $c_{u,v}=\frac{(u+v)!}{u!\,v!}$, $0\leq u,v,u+v\leq p-1$. Donc $c_{ij}\neq0\ (p)\iff c_{\unsure{v|r},s}\neq0\ (p)$.

Il reste à voir que si $w\geq p$, \deleted{soit $w=p+w'$} \inserted{alors} $c_{i,j}\equiv0\ (p)$. Soit $w^{\unsure{1|\prime|{}}}=p+w'\quad 0\leq w'\leq p-2$, donc $i+j=(r+s+1)p+w'=\underbrace{(t+1)}_{t'}p+w'$,
\[
(X+Y)^{i+j}=(X^p+Y^p)^{t'\deleted{\unsure{r}}}(X+Y)^{w'}=\sum_{\substack{0\leq\alpha\leq w'\\0\leq\beta\leq t'}}\binom{w'}{\alpha}\binom{t'}{\beta}X^{\alpha+p\beta}Y^{(w'-\alpha)+p(t'-\beta)}.
\]
\unsure{Toujours un|Trouver un|Toujours est|Terme en} $X^iY^j$ \unsure{correspond|correspondant} à $\alpha=u,\beta=r$, il faudrait $u=\alpha\leq w'=u+v-p$ donc $v\geq p$, ce qui n'est pas, donc il n'y a pas de
% archive sheet number: 19

